\GalSourcePageBoundary{088}{L0087}{59}{59}

J'observe maintenant que ce groupe est non primitif, en sorte
que toutes les lettres où le rapport des deux indices est le même
sont des lettres conjointes. Si l'on remplace par une seule lettre
chaque système de lettres conjointes, on aura un groupe dont
toutes les substitutions seront de la forme
\[
\biggl(b_{\frac{k_{1}}{k_{2}}},\
      b_{\frac{m_{1} k_{1} + n_{1} k_{2}}{m_{2} k_{1} + n_{2} k_{2}}}\biggr),
\]
$\dfrac{k_{1}}{k_{2}}$ étant les nouveaux indices. En remplaçant ce rapport par un
seul indice~$k$, on voit que les $p + 1$~lettres seront
\[
b_{0}, \quad b_{1}, \quad b_{2}, \quad b_{3}, \quad \dots, \quad b_{p-1}, \quad b_{\frac{1}{0}},
\]
et les substitutions seront de la forme
\[
\biggl(k,\ \frac{mk + n}{rk + s}\biggr).
\]

Cherchons combien de lettres, dans chacune de ces substitutions,
restent à la même place; il faut pour cela résoudre l'équation
\[
\GalPrintedError{(rk + s)k - m(mk + n) = 0}{W10-PE006},
\]
qui aura deux, ou une, ou aucune racine, suivant que $(m - s)^{2} + 4nr$
sera résidu quadratique, nul ou non résidu quadratique. Suivant
ces trois cas, la substitution sera de l'ordre $p - 1$, ou~$p$, ou~$p + 1$.

On peut prendre pour type des deux premiers cas les substitutions
de la forme
\[
(k,\ mk + n),
\]
où la seule lettre~$b_{\frac{1}{0}}$ ne varie pas, et de là on voit que le nombre
total des substitutions du groupe réduit est
\[
(p + 1)p(p - 1).
\]

C'est après avoir ainsi réduit ce groupe que nous allons le
traiter généralement. Nous chercherons d'abord dans quel cas un
diviseur de ce groupe, qui contiendrait des substitutions de
l'ordre~$p$, pourrait appartenir à une équation soluble par radicaux.

Dans ce cas, l'équation serait primitive et elle ne pourrait être
